% Section: Ukrainian abstract (анотація)
% =====================================================================
\begin{uabstract}
{\small\noindent Послідовні рекомендаційні системи на основі самоуваги ігнорують
насичений контрольований сигнал: інтенсивність кожної взаємодії
($5$-зірковий рейтинг, кількість годин гри, час перебування на
сторінці, вартість покупки). Ми висунули гіпотезу, що введення
інтенсивності в обчислення уваги повинно покращити точність top-$k$
пропорційно до того, наскільки інформативним є базовий сигнал. Ми
перевіряємо гіпотезу за допомогою трьох модифікацій SASRec
(IA-SASRec-Add, IA-SASRec-Mul, IA-SASRec-Val), що відповідають
трьом алгебраїчним позиціям, де інтенсивність може бути введена у
вираз уваги. Кожен варіант додає один скаляр $\lambda$, який
оптимізується під час навчання, на кожен шар уваги, спроектований
так, щоб при $\lambda{=}0$ точно відтворювався базовий SASRec:
модель є строгою надмножиною базової моделі. Ми проводимо оцінку на
семи наборах даних (MovieLens$\times2$, Amazon$\times2$,
Steam$\times3$) з п'ятьма ініціалізаціями на конфігурацію, за
протоколом значущості з двома критеріями та перевіркою стабільності
(попарний $t$-критерій на рівні ініціалізації та попарний бутстрап
на рівні користувача з $B{=}2{,}000$, захищений вимогою
мажоритарної згоди знаків довірчих інтервалів бутстрапу по
ініціалізаціях). Згідно з цим протоколом жоден варіант не досягає
стійкого покращення за метриками NDCG@10, Recall@10 або MRR@10 на
жодному з наборів. Отримана картина перевертає гіпотезу: саме на
наборі даних Steam (з найбільш насиченим сигналом інтенсивності)
спостерігаються найбільші статистично значущі погіршення результатів
(до $-9{,}5\%$ за показником NDCG@10, $p<0{,}05$). Аналіз механізму
на основі навченого параметра $\lambda$ пояснює причину невдачі:
$\lambda$ вільно адаптується на наборах даних, де інтенсивність не є
інформативною (до $0{,}07$ на ML-1M), але залишається на рівні
$0{,}55$--$0{,}99$ на Steam, де ці модифікації демонструють найгіршу
якість роботи. Архітектурний ``вимикач при $\lambda{=}0$'' є
принципово доступним, але на практиці не реалізований. Ми
стверджуємо, що один скалярний вентиль на кожен шар є занадто
обмеженим засобом контролю для інжекції інтенсивності, і пропонуємо
три конкретні архітектурні рішення, спрямовані на подолання цього
обмеження.

\par\smallskip
\noindent\textbf{Ключові слова:} послідовні рекомендації;
самоувага; SASRec; інтеграція додаткової інформації; увага з
урахуванням інтенсивності; відтворюваність; негативний результат;
аналіз механізму.}
\end{uabstract}
